% !TeX root = ../main.tex
% ---------------------------------------------------------------------------
% Squeezing, Swelling and High-Stress Ground
% Scope extension: expands the large-deformation material in Chapter 6.
% Key sources: Hoek & Marinos (2000); Barla on squeezing; Kaiser et al. on
% rockburst support; Martin et al. (1999) depth of spalling; Cai & Kaiser,
% Rockburst Support Reference Book.
% ---------------------------------------------------------------------------
\chapter{Squeezing, Swelling and High-Stress Ground}
\label{ch:squeezing-rockburst}


\section{Classification of difficult ground behaviour}
\label{sec:difficult-ground}
\scopenote{A behaviour-based taxonomy: squeezing, swelling, spalling,
bursting, ravelling and flowing, with the strength-to-stress ratio and
structure that produce each.}

\section{Squeezing ground}
\label{sec:squeezing-ground}

\subsection{Squeezing: mechanism and prediction}
\label{sec:squeezing-prediction}
\scopenote{The Hoek \& Marinos strain chart and $\scm/\pzero$ criterion,
Singh's and Jethwa's empirical criteria, and the expected convergence
percentage bands.}

\subsection{Time-dependent deformation and creep}
\label{sec:creep}
\scopenote{Distinguishing face-advance deformation from genuinely
time-dependent creep, viscoplastic models, and the consequences for support
timing and final lining loading.}

\subsection{Support strategies for squeezing ground}
\label{sec:squeezing-strategies}
\scopenote{Yield-and-control with deformation slots and sliding sets;
resist-and-stiffen with heavy support; and avoidance through alignment or
excavation-method change. Cost and risk of each.}

\section{Deformation allowance and swelling}
\label{sec:deformation-allowance-and-swelling}

\subsection{Over-excavation and deformation allowance}
\label{sec:over-excavation}
\scopenote{Designing an allowance for convergence so that the finished profile
is achieved; reprofiling as a programme risk; and the interaction with lining
thickness.}

\subsection{Swelling rock and soil}
\label{sec:swelling}
\scopenote{Anhydrite-to-gypsum conversion and smectite hydration; swelling
pressure testing; invert heave and the design responses of resistance versus
accommodation.}

\section{Brittle failure at high stress}
\label{sec:brittle-failure-at-high-stress}

\subsection{High-stress brittle failure}
\label{sec:brittle-failure}
\scopenote{Spalling, slabbing and notch formation; the cohesion-weakening
friction-strengthening behaviour that Mohr--Coulomb and Hoek--Brown fail to
capture near the boundary.}

\subsection{Depth of failure prediction}
\label{sec:depth-of-failure}
\scopenote{The Martin et al.\ linear relation between depth of failure and the
maximum boundary stress ratio, and its use in sizing bolt length and
retention.}

\section{Energy and excavation stability}
\label{sec:energy-and-excavation-stability}

\subsection{Energy changes on excavation}
\label{sec:energy-changes}
\scopenote{Strain energy stored in the rock before excavation, the energy
released when the opening is created, the work done on the support, and the
excess energy available to drive violent failure. After Brady \& Brown
Ch.~10.}

\subsection{Energy release rate and excavation stability}
\label{sec:energy-release-rate}
\scopenote{Energy release rate as a design index for deep excavation, its use
in sequencing and layout, and the stiffness relationship between the loading
system and the failing volume that determines whether failure is stable or
violent.}

\section{Rockburst and dynamic support}
\label{sec:rockburst-and-dynamic-support}

\subsection{Rockburst types and mechanisms}
\label{sec:rockburst-types}
\scopenote{Strainburst, pillar burst and fault-slip burst; energy release and
seismic source characterisation; the difference between a stress problem and a
seismic problem.}

\subsection{Dynamic support design}
\label{sec:dynamic-support}
\scopenote{Energy demand from ejection velocity and thickness, energy capacity
of bolts, mesh and shotcrete, and the requirement for system compatibility
rather than component capacity.}

\subsection{Destressing and preconditioning}
\label{sec:destressing}
\scopenote{Destress blasting and hydraulic fracturing to move stress away from
the boundary; evidence for effectiveness and the operational cost.}

\section{Monitoring and case studies}
\label{sec:monitoring-and-case-studies}

\subsection{Monitoring in high-stress and squeezing ground}
\label{sec:high-stress-monitoring}
\scopenote{Convergence arrays, extensometers, microseismic monitoring, and
their integration with excavation control.}

\subsection{Case studies}
\label{sec:squeezing-cases}
\scopenote{Gotthard Base Tunnel squeezing zones, deep hard rock mining
seismicity, and deep hydropower headrace tunnels.}
